\documentclass[10pt,a4paper]{amsart}
\usepackage[margin=19mm]{geometry}
\usepackage{amsmath,amssymb,mathtools}
\usepackage[scr=rsfs,cal=euler]{mathalfa}
\usepackage{xfrac}
\usepackage[unicode,hidelinks,bookmarks=false]{hyperref}
\newcommand{\ie}{i.e.\ }
\newcommand{\cf}{cf.\ }
\newcommand{\resp}{respectively\ }
\newcommand{\Subsection}{\subsection}
\providecommand{\nobreakdash}{\nobreak\hbox{-}}
\setlength{\emergencystretch}{2em}
\multlinegap0pt
\usepackage{amssymb}
\usepackage[all]{xy}
\usepackage{multirow}
\usepackage{booktabs}
\usepackage[shortlabels]{enumitem}
\usepackage{natbib}
\newtheorem{lemma}{Lemma}
\def\bg{\mathbf{g}}
\def\hil{\mcl{H}}
\def\mcl#1{\mathcal{#1}}
\def\sbracket#1{\langle #1\rangle}
\title{Why High-rank Neural Networks Generalize?: An Algebraic Framework with RKHSs}
\author{Yuka Hashimoto, Sho Sonoda, Isao Ishikawa, Masahiro Ikeda}
\date{Source: arXiv:2509.21895v2 (CC BY 4.0)}
\begin{document}
\maketitle

\noindent\emph{Source and licence.} Why High-rank Neural Networks Generalize?: An Algebraic Framework with RKHSs. Version arXiv:2509.21895v2 (CC BY 4.0), distributed by arXiv under the Creative Commons Attribution 4.0 International licence, http://creativecommons.org/licenses/by/4.0/, as recorded in the arXiv licence metadata for this exact version and shown on the abstract page https://arxiv.org/abs/2509.21895v2. Full licence notice: \texttt{licenses/arxiv-2509.21895-CC-BY-4.0.txt}.\par\smallskip

\noindent\textit{Editorial scope.} Counted unit: the article's lemma lem:iota\_surjective (source0.tex lines 838--840). The article's complete original proof of it is delivered with it (source0.tex lines 841--857, one \texttt{proof} environment of 17 lines). No mathematics is written, repaired or re-derived: the statement and the proof are the article's own lines, in the article's own notation, and no retained line is substituted or re-flowed.  No further source-local context is needed: every cross-reference of every retained line resolves inside the two delivered spans.  Nothing was omitted from any retained span, so the package carries no omission record.  No retained line contains a citation, so the wrapper carries no bibliography and no bibliography entry.  No retained line contains a figure environment, an image-inclusion command or drawing code, so no image is delivered and the package declares no asset dependency.  Editorial formatting only, in addition: an AMS \texttt{amsart} preamble that restates the article's own declarations and macros used by the retained lines, namely the environments \texttt{lemma} on separate counters, together with the standard packages those lines need.  The retained environments print in this document as Lemma 1 in order of appearance, whereas the article prints the same items with the numbering of its own full document.  The complete native source of the article as distributed in the arXiv e-print for this exact version is bundled unchanged as \texttt{sources/185776/source0.tex}.
\begin{lemma}\label{lem:iota_surjective}
The map $\iota$ preserves the norm and is surjective.
\end{lemma}

\begin{proof}
By definition, $\iota$ is a linear map that maps $\tilde{\phi}(\bg)\in\mcl{K}_0$ to $\phi(\bg)\in\mcl{R}_{k,0}$.
Thus, we have $\iota(\mcl{K}_0)=\mcl{R}_{k,0}$.

For $h\in\mcl{K}_0$, there exist $n\in\mathbb{N}$, $\bg_1,\ldots,\bg_n\in G^L$, and $c_1,\ldots,c_n\in\mathbb{C}$ such that $h=\sum_{i=1}^nc_i\tilde{\phi}(\bg_i)$.
We have
\begin{align*}
\Vert\iota(h)\Vert^2_{\mcl{R}_k} = \bigg\Vert\sum_{i=1}^n c_i \phi(\bg_i)\bigg\Vert^2_{\mathcal{R}_k} = \sum_{i,j=1}^n\overline{c_i}c_jk(\bg_i,\bg_j)
=\sum_{i,j=1}^n\overline{c_i}c_j\sbracket{\tilde{\phi}(\bg_i),\tilde{\phi}(\bg_j)}_{\hil}
=\Vert h\Vert^2_{\hil}.
\end{align*}
Thus, $\iota$ preserves the norm, and in particular, it is bounded.

For any $r\in \mcl{R}_k$, there exists a sequence $r_1,r_2,\ldots \in \mcl{R}_{k,0}$ such that $\lim_{i\to\infty}r_i=r$.
Since $\iota(\mcl{K}_0)=\mcl{R}_{k,0}$, there exists $h_i\in\mcl{K}_0$ such that $\iota(h_i)=r_i$ for $i=1,2,\ldots$.
Thus, we have $r=\lim_{i\to\infty}r_i=\lim_{i\to\infty}\iota(h_i)=\iota(\lim_{i\to\infty}h_i)=\iota(h)$.
\end{proof}

\end{document}
